\documentclass[10pt,a4paper]{article}
\usepackage[margin=23mm]{geometry}
\usepackage{amsmath,amssymb,booktabs,array,longtable}
\usepackage[hidelinks]{hyperref}
\setlength{\parindent}{0pt}
\setlength{\parskip}{4pt}
\setlength{\emergencystretch}{2em}
\newcommand{\dd}{\mathrm d}
\newcommand{\pos}[1]{[#1]_+}
\title{Urban effective-service control:\\a proposed timestep methodology}
\author{Selective-discreteness urban review}
\date{27 October 2026}
\begin{document}
\maketitle

\textbf{Feedback 10 status.}
This document preserves an earlier candidate transcription. The user selects
shared left/through permission for the starting bay case, actual-admission
commitment, unassigned movement-priority residuals and average receiving A.
The logical allocator may inspect other movements after a failed allocation;
its physical access/interface with the common feeder remains open.

Frozen waiting targets and old-target rationing are unselected conditional
material. Q4 is withdrawn. A still has sequential head-clearance and
current-head-sharing dependencies. The updated coupled model is not yet
defined by the old equations below. This is an answer/status update only;
see the separate \href{urban_feedback10_response.md}{feedback 10 answer}.

\section{The traffic question and admitted world}

A vacant left-turn bay is useful only if a turner can reach it. A through
block in the common feeder can prevent access: additional left service need
not help, while through service can remove the obstacle. The question is
whether continuous mass with exact local semantic order predicts useful
control responses that sound aggregate explanations miss.

Retain the accepted isolated four-leg task, V0--V4, fixed routes and four
protected effective-service groups: NS through/right, NS left, EW through/right,
EW left. Let
\begin{equation}
 g_p\geq0,\quad \sum_{p=1}^4g_p=1,\qquad
 \mu_{\ell m}(g)=\min\{C_\ell,s_{\ell m}g_{p(m)}\}.
 \label{eq:action-new}
\end{equation}
Fractions allocate average service, without defining executable phase
intervals. Perturbations preserve the sum and nonnegativity. Saturation $s$
and section capacity $C$ have units veh-eq/s. No phase waveform, microscopic
lane change, route choice or network-wide lane expansion is introduced.

V0 uses independent movement-exclusive paths. A through/right service class
is order-neutral only if its members face equivalent receiving restrictions,
or these are nonbinding. V1 has a shared stop-line lane. V2/V3 have the same
finite-bay geometry: V2 is constructed so order does not restrict access;
V3 retains access-sensitive order. V4 varies declared geometry, demand and
receiving conditions on held-out cases. A V2 label alone does not prove an
order-neutral trajectory.

Physically split an approach into upstream aggregate section $U$, common
local feeder $F$, and parallel downstream lanes $b$ beginning at an internal
fork. Bay and through continuation have the same fork-to-stop-line distance
in the symmetric core, with separately specified capacity/storage. Path
length is $L_U+L_F+L_b$; these are disjoint storage owners. V1 omits the fork;
V0 uses independent paths, without a synthetic common FIFO feeder. This
internal-bay composition extends neutral M2; it is not implemented behavior.

Permissions are an explicit structural input, not inferred from lane count.
The exclusive-target illustration used for the matched closure below is
$L\to B$ and $T/R\to T$. It is one candidate, not the selected primary
matrix. The expanded decision appendix in the compact report also considers
left turns in both bay and shared continuation; feedback 10 selects that
configuration for the starting bay scenario. It requires a
matching comparator, rather than silently applying the exclusive-target
ablation. Both V2 and V3 must share the chosen permissions.

\section{State, timing and exact order}

Each physical section $e$ owns cumulative counts $N^+_{e,k},N^-_{e,k}$,
occupancy $n_{e,k}=N^+_{e,k}-N^-_{e,k}$ and propagation history. A source
owns waiting mass $Q_{o,k}$. Ordered ledgers retain route-table reference,
route position, continuous mass and eligibility provenance. Route position
advances at the actual junction crossing, not at internal $U/F/b$ crossings.

\textbf{Historical state qualification.}
The remaining-target vector described next belongs to the earlier waiting
proposal. New unadmitted mass has no such binding vector under feedback 10;
realized lane occupancy/order remain fixed.

A canonical homogeneous cohort has mass $q$, one current movement and
route reference/position, one owner and compatible contiguous timing. A
remaining target vector $a_b\geq0$, $\sum_ba_b=q$, is a view of mass still
upstream, not another physical queue. One cohort may split among legal lanes
without an artificial order between its children. Different cohorts are
processed in source order; every realized lane retains that subsequence.

Storage fragments do not define traffic behavior. Before allocation/admission,
coalesce exactly compatible adjacent fragments into maximal cohorts.
Compatibility requires identical route/position, owner, permissions, affine
eligibility and, for assigned cohorts, target proportions. Retain M1's exact
affine-line test, without tolerance-based mixing. Different commitments or
timing are not identical metadata. This canonicalization prevents an
arbitrary storage split changing the homogeneous group seen by allocation.
It is a proposed transcription requirement, not a new integer packet unit.

The grid is a movement view, not the order authority. If ledger entry $r$
occupies mass interval $I_r$ and cell $c$ is $B_c$,
\begin{equation}
 M_{cm}=\sum_{r:m(r)=m}|I_r\cap B_c|.
 \label{eq:grid-new}
\end{equation}
Discharge consumes the ledger, not an inferred order from $M$. Rank is mass,
not metres. Grid resolution cannot change the exact-ledger traffic law.

\section{Inherited boundary-history propagation}

Use $t_k=k\Delta t$, empty prehistory and triangular-FD sections with
delays $\tau^f_e,\tau^b_e$, capacity $C_e$ and storage
$K_e=C_e(\tau^f_e+\tau^b_e)$. Require
$\Delta t\leq\min_e(\tau^f_e,\tau^b_e)$; a short-section violation requires
reconsidering discretization, not clipping delays.
For $\theta=(t-t_h)/\Delta t$,
\begin{equation}
 {\cal H}^{\pm}_e(t)=(1-\theta)N^\pm_{e,h}+\theta N^\pm_{e,h+1},
 \quad t_h\leq t\leq t_{h+1}.
 \label{eq:interpolation-new}
\end{equation}
Histories are zero for $t<0$. Compute pre-update allowances:
\begin{align}
 A_{e,k}&=\pos{{\cal H}^+_e(t_{k+1}-\tau^f_e)-N^-_{e,k}},\nonumber\\
 S_{e,k}&=\min(C_e\Delta t,A_{e,k}),\label{eq:sending-new}\\
 R_{e,k}&=\min\{C_e\Delta t,
 \pos{K_e+{\cal H}^-_e(t_{k+1}-\tau^b_e)-N^+_{e,k}}\}.
 \label{eq:receiving-new}
\end{align}
These are masses per update. The first $S_{e,k}$ resident ledger mass is the
sending tranche; the first $A_{e,k}$ is available for logical assignment.
This inherits discrete readiness; it does not assert that all the tranche
was physically ready at every instant within the step. New admissions
cannot receive same-section service in this update.

For actual entrance $b_{e,k}$ and exit $y_{e,k}$,
\begin{equation}
 N^+_{e,k+1}=N^+_{e,k}+b_{e,k},\quad
 N^-_{e,k+1}=N^-_{e,k}+y_{e,k},\quad
 n_{e,k+1}=n_{e,k}+b_{e,k}-y_{e,k}.
 \label{eq:state-new}
\end{equation}
Today's departure creates no immediate extra receiving room. Each internal
transfer appears once with each sign in its two physical owners.

After selecting flows, reconstruct each receiving lane's entrance segments
from its boundary inflow and retained child order. For total inflow $b>0$,
a child occupying entrance rank $[r,r+q']$ receives time segment
$[t_k+\Delta t\,r/b,t_k+\Delta t(r+q')/b]$, with the receiving free-flow
shift. This applies M1-style affine reconstruction to the new section graph.
Do not retain internal service-event times as new propagation knots.
Cross-input ties on outgoing links are disclosed bookkeeping; this isolated
task has no later controlled junction where such a tie could confer an
order advantage.

\section{Known capability inside allocation}

This is the proposed exact completion of feedback 9. Separate feasible entry
now from remaining intention. Do not first fix a target without known limits.

For this joined proposal, first new assignment occurs at the fork's
eligible front, before physical admission is guaranteed. The older prototype
placed assignment at entry to $F$; that location was not selected in the
principle discussion. This proposed location change is explicit: it makes
current branch limits available without forecasting travel through $F$.
Existing valid waiting targets are still preserved.

At the fork process canonical cohorts in physical source order. Lane-load
estimate $h_b$ is pre-update resident branch mass plus target mass assigned
by earlier cohorts processed this update, including pending mass. Later
cohorts therefore see earlier assignments. Do not subtract predicted signal
discharge to invent service-aware driver knowledge. Compare raw mass between
matched legal lanes; other perception/cost models remain outside this draft.

Let $s$ be remaining common sending, $u_b$ remaining receiving on reachable
legal branches, and $q$ the available cohort mass. For a fresh cohort:
\begin{align}
 V&=\min\{q,s,\sum_bu_b\},\nonumber\\
 x&=\underset{0\leq x_b\leq u_b,\ \sum_bx_b=V}{\arg\min}
                    \frac12\sum_b(h_b+x_b)^2,\nonumber\\
 x_b&=\min\{u_b,\pos{\lambda-h_b}\},\quad \sum_bx_b=V.
 \label{eq:feasible-allocation}
\end{align}
The strictly convex objective selects unique $x$; the threshold itself need
not be unique on a fully capped plateau. For $V=0$ set $x=0$. This is
capped continuous queue filling for one cohort, not a heterogeneous equilibrium.

\textbf{Unselected after feedback 10.}
The following waiting-target and old-target formulas are preserved for
traceability, not adopted as the current model. Their premise is absent
under commitment only at actual admission.

To retain logical assignment of waiting traffic, assign remaining intention
only after feasible entry has been selected:
\begin{equation}
 z_b=\pos{\nu-h_b-x_b},\quad \sum_bz_b=q-V,\qquad a_b=x_b+z_b.
 \label{eq:pending-allocation}
\end{equation}
Use only legal eventual targets. No current receiving is consumed by $z$,
and it promises no current physical access. The new target $a$ then stays
fixed. This second queue fill is an explicit proposed waiting-intent
convention, not a user-approved unconstrained-first allocation law.

For valid old remaining targets $a$, do not allocate again:
\begin{equation}
 \bar u_b=\min(u_b,a_b),\quad V=\min\{q,s,\sum_b\bar u_b\},\quad
 x_b=\min\{\bar u_b,\pos{\lambda-h_b}\},\quad z_b=a_b-x_b,
 \label{eq:old-commitment}
\end{equation}
with $\sum_bx_b=V$. The same proposed queue criterion selects feasible
realization within fixed targets. It does not retarget when capacity changes;
nor does it claim feedback 9 selected former Candidate B.

\textbf{Scope of the old prefix gate.}
The physical prefix statement below applies to an occupied single common
lane. It is not the selected general rule for logical allocation-server
processing. Reconciling movement-specific inspection with the physical
interface remains open in this answer-only update.

Append $x_b$ to branch $b$, reduce sending/receiving by actual entry, and
set $h\leftarrow h+a$. Count $a=x+z$ once, not $x+a$. The residual $q-V$
and targets $z$ remain with the upstream owner. On a genuine common feeder,
any residual older mass closes the physical prefix gate for later cohorts
this update. Later logical assignment can continue with $V=0$; it cannot
authorize bypass. Independent paths have independent gates.

Cut an unassigned cohort at the inherited eligibility boundary before
creating targets, preserving the not-yet-eligible tail. Under empty loading
and assignment at this fork, assigned mass was already eligible here;
its residual remains eligible while waiting at the same owner. A sending
quota can still restrict its actual entry through $V$. No extra
proportional retargeting of a newly ready part is required. Within a
homogeneous admitted prefix, lane amounts split that prefix; one target
child is not artificially placed before another.

\textbf{Why the dependency matters.}
For $h=(0.2,0)$, $q=1$, $u=(0.3,0.5)$ and $s=0.6$,
$x=(0.2,0.4)$, $z=(0.2,0.2)$, $a=(0.4,0.6)$:
$0.4$ remains upstream. If $u=(0,1)$ and $s=1$ instead,
$x=a=(0,1)$. A target $(0.4,0.6)$ frozen without those known limits would
admit only $0.6$. Conversely, if that vector really is an older valid target,
the $0.4$ cannot silently be reassigned.

\section{Elapsed service and current-head receiving reuse}

Expose only each lane's current ordered head in its pre-state sending
tranche, with rate cap $c_\ell=\mu_{\ell m}(g)$. A zero-rate head blocks that
lane. For the proposed timestep node solver define
\begin{equation}
 \rho_j=R_{j,k}/\Delta t.
 \label{eq:interval-rate}
\end{equation}
This is an explicit step-average computational convention, not the physical
availability curve inside the interval. It spends at most the inherited
receiving allowance and writes endpoint history only. Event histories remain
deferred.

For outgoing link $j$, collect current heads from leg $i$. In this isolated
topology each leg-to-output movement has one relevant service group:
\begin{align}
 C_{ij}&=\sum_{\ell\in{\cal L}_{ij}}c_\ell,\quad
 w_{ij}=g_{p(i,j)},\quad V_j=\min(\rho_j,\sum_iC_{ij}),\nonumber\\
 r_{ij}&=\min(C_{ij},\theta_jw_{ij}),\quad \sum_i r_{ij}=V_j,
 \label{eq:leg-rates-new}\\
 r_\ell&=\min(c_\ell,\eta_{ij}),\quad
 \sum_{\ell\in{\cal L}_{ij}}r_\ell=r_{ij}.
 \label{eq:lane-rates-new}
\end{align}
Remove zero-cap competitors; $V_j=0$ gives zero rates. Weight each leg once,
with capped reuse; share equally among its lanes subject to individual caps.
Independent outputs have independent solves because each lane has one
current head/output.

Hold rates until the earliest head/tranche exhaustion or step end:
\begin{equation}
 \delta=\min\left\{\Delta t-\sigma,\
 \min_{\ell:r_\ell>0}\frac{\min(q_\ell,S_\ell^{\rm rem})}{r_\ell}\right\}.
 \label{eq:head-step-new}
\end{equation}
Advance every active lane by $r_\ell\delta$, process tied events together,
expose new heads and recompute sharing. Stop if all rates vanish. A finite
canonical ledger gives finitely many head events; use no mass cutoff or
software scan priority. This internal clock calculates one step's output;
it creates no new observation grid.

For ready left mass $q_L$ followed by sufficient through mass, with
nonbinding receiving/capacity,
\begin{equation}
 y_L=q_L,\qquad y_T=s_Tg_T
       \left(\Delta t-\frac{q_L}{s_Lg_L}\right),
 \quad 0<\frac{q_L}{s_Lg_L}<\Delta t.
 \label{eq:order-response-new}
\end{equation}
With $\Delta t=s_L=s_T=1$, $q_L=0.2$ and $(g_L,g_T)=(0.4,0.6)$,
$y_T=0.3$. Transferring small $\epsilon$ from through to left gives
$\dd y_T/\dd\epsilon=0.25$: faster head clearance outweighs the smaller
through rate on this branch. This hand case is not a complete-TST or
usefulness result. Other active constraints can change or nullify it.

\section{Complete update and feasibility}

Read pre-state histories/ledgers; compute all S/R and logical assignment
fronts; process source/internal admissions and forks; solve stop-line service
on pre-state tranches; commit paired flows; reconstruct boundary eligibility
and grid views; update histories, sources, pending views and sinks.
Independent calculations may be evaluated in either software order, but
no mechanism spends another's allowance or serves new same-section inflow.

A one-in/one-out internal connection transfers
$f=\min(S_{\rm up},R_{\rm down})$ from its resident prefix.
Release $D_{o,k}$ joins source waiting at the boundary; admission is bounded
by $Q_{o,k}+D_{o,k}$ and first-section receiving, in source order. Independent
source paths own separate queues. A sink removes $\min(S,\bar S_k)$,
where $\bar S_k$ is its declared allowance. The central junction uses the
current-head law above; there is no hidden additional conflict rule.

The timestep restriction gives $S_{e,k}\leq n_{e,k}$ and
$R_{e,k}\leq K_e-n_{e,k}$. Thus
\begin{equation}
 0\leq n_{e,k}+b_{e,k}-y_{e,k}\leq K_e,\qquad
 0\leq b_{e,k}\leq R_{e,k},\quad0\leq y_{e,k}\leq S_{e,k}.
 \label{eq:storage-proof-new}
\end{equation}
Internal paired flows and route-preserving splits cancel, giving
\begin{equation}
 \sum_oQ_{o,k}+\sum_en_{e,k}+Z_k=D_{<k},
 \label{eq:population-new}
\end{equation}
with cumulative sink exit $Z_k$ and exogenous release $D_{<k}$.
This holds route by route; pending commitments add no stock.
These are definition-level arguments, not tested implementation evidence.

\section{One-step derivatives}

Work on a fixed order/structural and active-constraint signature. Differentiate
continuous split masses, queues, pending targets and event durations, without
differentiating route labels or permissions.

For capped queue fill with free set $ {\cal A}$ and upper-capped set
$ {\cal B}$, if $ {\cal A}\ne\varnothing$,
\begin{align}
 \lambda&=\frac{V-\sum_{{\cal B}}u_b+\sum_{{\cal A}}h_b}{|{\cal A}|},
 &\dot\lambda&=\frac{\dot V-\sum_{{\cal B}}\dot u_b+
                       \sum_{{\cal A}}\dot h_b}{|{\cal A}|},\nonumber\\
 \dot x_b&=\dot\lambda-\dot h_b\ (b\in{\cal A}),&
 \dot x_b&=\dot u_b\ (b\in{\cal B}),\quad
 \dot x_b=0\ (\text{zero set}).
 \label{eq:allocation-jacobian-new}
\end{align}
Use $\bar u$ for old targets, retaining $\dot a$ on active target caps.
Differentiate $V$ through its active minimum. Fresh waiting intent uses
the same free-set calculation with base $h+x$, total $q-V$, no finite
upper caps; old intent has $\dot z=\dot a-\dot x$.
On fully capped branches $\dot x=\dot u$; on a stable zero-volume branch
$\dot x=0$. Transitions may be kinks, not ordinary derivatives.

For leg sharing with uncapped set $ {\cal U}$ and capped set $ {\cal B}$,
\begin{equation}
 \dot\theta=
 \frac{\dot V-\sum_{{\cal B}}\dot C_{ij}
                -\theta\sum_{{\cal U}}\dot w_{ij}}
      {\sum_{{\cal U}}w_{ij}},\qquad
 \dot r_{ij}=
 \begin{cases}\dot C_{ij},&i\in{\cal B},\\
 w_{ij}\dot\theta+\theta\dot w_{ij},&i\in{\cal U}.
 \end{cases}
 \label{eq:priority-jacobian-new}
\end{equation}
Within-leg sharing uses the unweighted analogue. For an interior head event
$\delta=q/r$, $\dot\delta=\dot q/r-q\dot r/r^2$.
For total output $y=\sum_a r_a\delta_a$,
\begin{equation}
 \dot y=\sum_a(\dot r_a\delta_a+r_a\dot\delta_a).
 \label{eq:service-jacobian-new}
\end{equation}
Do not freeze head-clearance times. The final interval's derivative is minus
the sum of previous duration derivatives. Ties need compatible branches or
explicit one-sided analysis.

Differentiate propagation and conservation through their active branches.
Fixed delayed query times have fixed interpolation weights, but the history
values retain sensitivity. A cut within an affine eligibility segment
$[a,b]$ uses fraction $(t-a)/(b-a)$ and its full derivative. Point segments,
zero mass and changes in ledger topology require their boundary branches;
the formula is not a global smoothness claim.

\section{From discharge to complete TST}

Collect all retained numerical state in $X_k$ and write the defined update as
$X_{k+1}=F_k(X_k,u_k)$. The inherited left-interval objective is
\begin{align}
 J_{\Delta t}&=\Delta t\sum_{k=0}^{H-1}
       \left[\sum_o(Q_{o,k}+D_{o,k})+\sum_en_{e,k}\right],
 \label{eq:tst-new}\\
 P_0&=0,\quad P_{k+1}=F_{k,X}P_k+F_{k,u}\dot u_k,\quad
 \dot J=\sum_k(c_{k,X}P_k+c_{k,u}\dot u_k).
 \label{eq:chain-new}
\end{align}
Within control $\dot u_k=d$; under candidate-independent continuation
$\dot u_k=0$, while $P_k$ keeps propagating.
Candidate-independent preparation gives $P_0=0$; otherwise differentiate it.
No per-vehicle trajectory derivative is needed.

At fixed state the direct partial of this step's cost with respect to service
is zero. Changed service can still alter next queues, access, spillback,
source waiting and eventual sink exit. Internal transfer, including a
stop-line crossing into another road, does not itself remove system population.
For fixed releases, an independent accounting identity is
\begin{equation}
 J_{\Delta t}=\Delta t\sum_{k=0}^{H-1}(D_{\leq k}-Z_k),\qquad
 \dot J_{\Delta t}=-\Delta t\sum_{k=0}^{H-1}\dot Z_k.
 \label{eq:tst-exit-new}
\end{equation}
This counts all affected mass without duplicate pending stock.

Propose a common integer horizon $H$ qualified to clear every compared
nearby feasible action, followed by zero occupancy. This chooses the earlier
common-horizon option, not a numerical horizon or universal clearance
theorem. Finite demand, acyclic routes, positive continuation service and
eventually available sinks support clearance. The chosen horizon must still
be qualified over the comparison neighborhood. A residual cutoff,
nonclearing case or changing integer stop index is not complete-TST
differentiability evidence.

\section{Fair and reconstructible aggregate comparisons}

Allow the original fine record through each rival's declared macro converter
or estimator. Disclose states, retained memory, observation times and
calibration data. Native subsets are legitimate; successful macro inference
counts against needing explicit order. A persistent ordered ledger is meso
information even inside a macro-labelled estimator. Separate retrospective
replay, observation-updated prediction and forecasting from a common prefix.
Each derivative claim also states which observed inputs are held fixed or
regenerated under the candidate action. A later observed outcome cannot
count as having been predicted before it was observed.

Add one controlled matched aggregate ablation, $B_{\rm mix}$, alongside
source-native B1/B2. It uses the same section graph, S/R, action, releases,
loading, continuation and objective, but retains route/movement totals per
section without an ordered ledger. Its following closure is defined here,
not borrowed wholesale from a native source.

At a common feeder use pre-state composition $\beta_m$. In the primary bay
geometry each movement has its declared branch:
\begin{equation}
 f_m=\beta_m V,\qquad
 V=\min\left\{S_F,\min_{m:\beta_m>0}\frac{R_{b(m)}}{\beta_m}\right\}.
 \label{eq:mix-fork-new}
\end{equation}
Ignore zero fractions; add new inflow after removing this mixture. A
homogeneous movement with multiple matched legal targets uses the same
constrained queue fill, splitting route counts proportionally within that
movement. Do not add permissions merely to create lane-choice freedom.

At a mixed stop-line lane set
\begin{equation}
 c_\ell=\min\left\{C_\ell,
 \left(\sum_{m:\beta_{\ell m}>0}
       \frac{\beta_{\ell m}}{s_{\ell m}g_{p(m)}}\right)^{-1}\right\}.
 \label{eq:mix-rate-new}
\end{equation}
A positive fraction at zero service makes $c_\ell=0$; pure lanes reduce to
the selective cap. This is homogeneous elapsed work, not recovered order.
Set positive priority
$w_\ell=\sum_m\beta_{\ell m}g_{p(m)}$ and output fractions $\beta_{\ell j}$.
For the primary topology, each leg has at most one competing lane per output.
Starting at zero, increase unfrozen $r_\ell$ proportionally to $w_\ell$
until a lane cap or output constraint
$\sum_\ell\beta_{\ell j}r_\ell\leq\rho_j$ binds. Freeze capped lanes and
all positively turning competitors of a bound output. Continue the remaining
lanes, processing simultaneous ties together, until all are fixed. This
defines a deterministic weighted progressive-fill rate law. Consume each
lane's pre-state sending quota, recomputing on quota exhaustion. Extending
this comparator to several same-leg competitors requires retaining the
agreed leg-level sharing hierarchy; no such extension is assumed here.

Neutral equivalence requires matched pure channels, or a single same-service
class with equivalent receiving and readiness. Under the declared primary
pure-channel cases both laws reduce to the same S/R/service minimum and
current receiving priorities. Equal mixed totals alone do not imply
equivalence. $B_{\rm mix}$ isolates a declared closure; it does not replace
strong native rivals or establish macro impossibility.
In particular, its full-FIFO mixture can restrict bay entry more than an
unblocked ordered prefix. A V2 failure is then evidence against that
aggregate closure, not proof that exact order is necessary. The principal
storage-aware rival must first pass the relevant unblocked/storage checks
before a V3 advantage can support the intended information claim.

\section{Native rivals and independent/external references}

B1's working preference is
\href{https://arxiv.org/html/2310.12249v2}{Wei et al.'s link queue model}: moving/queued
state with native effective green. Its complete causal link/node law,
geometry and initialization must be instantiated. The source's
short-free-flow-region limitation of its fixed-speed shortcut prevents
treating that shortcut as an unconditional spillback model.

B2's working preference is
\href{https://www.sciencedirect.com/science/article/pii/S0191261525000578}
{Gong et al.'s lane-exit allocation} anchor with a
sourced finite-storage/access adaptation. Exit allocation is not finite-bay
entry. \href{https://attap.umd.edu/wp-content/uploads/2022/05/TRC522.pdf}
{Liu--Chang} supplies a concrete lane-group storage/blockage source;
its native signal/timing and overflow variables need reconciliation before
a combined B2 is qualified. Native theorems do not automatically transfer.
These implementation-dependent choices stay pending under feedback 8;
no new class-selection vote is required.

R1 uses independent hand/reference cases: prefix obstruction, partial entry,
receiving reuse, rotation, route-mass balance and neutral limits. They are a
correctness floor, not external realism.

For R2 declare a static executable program with fixed cycle $C^{\rm cyc}$,
lost intervals $\ell_p$, $L=\sum_p\ell_p<C^{\rm cyc}$, order and offset. Propose
\begin{equation}
 G_p=(C^{\rm cyc}-L)g_p,\qquad
 s^{\rm eff}_{\ell m}=s^{\rm phys}_{\ell m}
                    (C^{\rm cyc}-L)/C^{\rm cyc}.
 \label{eq:sumo-map-new}
\end{equation}
Keep these settings, lane behavior and calibration fixed across model
directions and candidate actions. The mapping matches nominal average
opportunity, not startup loss or realized microscopic discharge.

Use a common finite route population and desired releases on declared
common boundaries. Integer reference vehicles are exogenous population,
not rounding of internal simulated mass. Every model receives that release
list or its disclosed conversion. SUMO's fixed lane behavior determines
realized occupancy; do not replay private selective assignments into only
one competitor.

External complete time is desired-release-to-arrival time: TripInfo
\texttt{duration} plus \texttt{departDelay}, without adding waiting time.
Track the master population so uninserted/unfinished demand cannot vanish
from a completed-trip statistic. Clear all demand or report noncompletion.
Common loading does not imply equal later states across simulators.
Separate state-conditioned replay from common-input prediction.
Test feasible perturbation signs/ranking externally before optimization.

\section{Status, evidence limits and source boundaries}

\textbf{Current qualification.}
The summary below describes the earlier candidate. Its frozen-waiting
construction is unselected after feedback 10, and the replacement physical
access/interface has not been derived in this response.


This proposal joins physical ownership, pre-state S/R, ordered allocation,
fixed waiting intent, prefix admission, elapsed service, endpoint history,
complete cost and conditional sensitivities. Newly explicit transcription
choices are: feasible-entry-first filling with subsequent waiting intention;
the same queue criterion within old targets; canonical cohorts; a
fork assignment location and explicit permission subfamily; a step-average
receiving clock; a qualified common horizon; and $B_{\rm mix}$.
They remain proposed definitions for review, not newly approved decisions.

Numerical geometry/FD parameters, timestep, grid resolution, horizon,
native B1/B2 qualification and R2 program values remain later work. No
simulator, AD/FD, optimization or SUMO experiment has run in this writing
pass. Reference correctness, invariants, neutral equivalence, meso-sensitive
forward behavior, derivative correctness and external usefulness need
separate evidence. Finite gradients or a falling internal objective cannot
establish useful control directions.

\begin{itemize}
\item M0 supplies repository LTM/history/TST; M1 supplies the route ledger,
prefix reconstruction and exact merge rule; M2 supplies ownership/grid
invariants. The approved M0 source basis is Yperman (2007) and UNsim
v0.12.0. This pass checks repository transcription, not fresh reproduction
of those complete originals.
\item The primary links beside B1/B2 identify their sources, not adoption
of their complete algorithms. Gong's indexed primary sections establish
scope; its complete algorithm is not independently reproduced here.
Liu--Chang's journal year is 2011; overflow notation needs verification.
\item Official SUMO
\href{https://sumo.dlr.de/docs/Simulation/Traffic_Lights.html}{signal-program}
and \href{https://sumo.dlr.de/docs/Simulation/Output/TripInfo.html}{TripInfo}
definitions.
They support interface/accounting definitions, not empirical equivalence.
\end{itemize}
\end{document}
